% Part: computability
% Chapter: recursive-functions
% Section: halting-problem

\documentclass[../../../include/open-logic-section]{subfiles}

\begin{document}

\olfileid{cmp}{rec}{hlt}
\olsection{د درېدنې مسئله}

په عمومي ډول د \emph{درېدنې مسئله} دا ده چې د يوې محاسبه کېدونکې
تابع د مشخصې~$e$ (لکه پروګرام) او يو عدد~$n$ په ورکړې سره پرېکړه وشي
چې د دې تابع محاسبه پر ننوت~$n$ درېږي، يعنې کومه پايله ورکوي، که نه۔
اېلن ټيورينګ په نوميالي ډول ثابته کړه چې دا مسئله پخپله د يوې محاسبه
کېدونکې تابع له لارې نۀ شي حل کېدے؛ يعنې تابع
\[
h(e, n) =
\begin{cases}
1 & \text{که محاسبه $e$ پر ننوت $n$ ودرېږي}\\
0 & \text{که داسې نۀ وي،}
\end{cases}
\]
محاسبه کېدونکې نۀ ده۔

د جزوي بازګشتي تابعو په چوکاټ کښې د يو پروګرام د مشخصې رول د کليني
د نورمال شکل په قضيه کښې ورکړل شوے شاخص~$e$ لوبولے شي۔ که $f$ جزوي
بازګشتي تابع وي، هر هغه~$e$ چې د نورمال شکل د قضيې معادله پوره کوي،
د~$f$ يو شاخص دے۔ د عدد~$e$ په ورکړې سره د نورمال شکل قضيه وايي چې
\[
\cfind{e}(x) \simeq U(\mu s \; T(e, x, s))
\]
جزوي بازګشتي ده، او د هرې جزوي بازګشتي $f\colon \Nat \to \Nat$
دپاره يو $e \in \Nat$ شته چې د ټولو~$x \in \Nat$ دپاره
$\cfind{e}(x) \simeq f(x)$ وي۔ په حقيقت کښې د هرې داسې~$f$ دپاره
يوازې يو نه، بلکې بې‌شمېره ډېر داسې~$e$ شته۔ \emph{د درېدنې تابع}~$h$
داسې تعريفېږي:
\[
h(e, x) =
\begin{cases}
1 & \text{که $\cfind{e}(x) \fdefined$ وي}\\
0 & \text{که داسې نۀ وي۔}
\end{cases}
\]
نو $h(e,x)=0$ هله او يوازې هله چې $\cfind{e}(x) \fundefined$ وي؛ د
نورمال شکل په ټاکلې شمېرنه کښې هر طبيعي عدد~$e$ د کومې جزوي بازګشتي
تابع شاخص دے، که هغه تابع په هر ننوت هم تعريف شوې نۀ وي۔
\textbf{د ژباړې د نسخې سمون (OLCMP-017):} سرچينې د نورمال شکل د
شمېرنې تر ټاکلو وروسته هم د داسې عدد بديل ساتلے و چې ګنې د هېڅ جزوي
بازګشتي تابع شاخص نۀ وي۔ خو پورته معادله هر طبيعي شاخص ته يوه جزوي
بازګشتي تابع ټاکي۔ دلته هغه ناممکن بديل له تشريح او قطري ثبوت څخه
لرې کړے شوے دے؛ د درېدنې د تابع تعريف او د قضيې دعوه هماغسې دي۔

\begin{thm}
\ollabel{thm:halting-problem}
د درېدنې تابع~$h$ جزوي بازګشتي نۀ ده۔
\end{thm}

\begin{proof}
که $h$ جزوي بازګشتي وای، دا تابع مو تعريفولے شوه:
\[
d(y) =
\begin{cases}
1 & \text{که $h(y, y) = 0$ وي}\\
\umin{x}{x \neq x} & \text{که داسې نۀ وي۔}
\end{cases}
\]
څرنګه چې هېڅ عدد~$x$ شرط $x \neq x$ نۀ پوره کوي،
$\umin{x}{x \neq x}$ نشته؛ نو $d(y) \fundefined$ هله او يوازې هله
چې $h(y,y) \neq 0$ وي۔ له دې تعريف څخه دا پايلې راځي:
\begin{enumerate}
\item د شاخص~$y$ دپاره، $d(y) \fdefined$ هله او يوازې هله چې
  $\cfind{y}(y) \fundefined$ وي۔
\item $d(y) \fundefined$ هله او يوازې هله چې $\cfind{y}(y)
  \fdefined$ وي۔
\end{enumerate}
که $h$ جزوي بازګشتي وای، $d$ به هم جزوي بازګشتي وه۔ نو د کليني د
نورمال شکل د قضيې له مخې دا يو شاخص~$e_d$ لري۔ د~$h(e_d,e_d)$ قيمت
وګورئ۔ دوه ممکن حالتونه شته، $0$ او~$1$۔
\begin{enumerate}
\item که $h(e_d, e_d) = 1$ وي، نو $\cfind{e_d}(e_d) \fdefined$ دے۔
  خو $\cfind{e_d} \simeq d$، او $d(e_d)$ هله تعريف شوے دے چې
  $h(e_d,e_d)=0$ وي۔ نو $h(e_d,e_d) \neq 1$۔
\item که $h(e_d, e_d) = 0$ وي، نو د شاخص~$e_d$ تابع
  $\cfind{e_d}(e_d) \fundefined$ ده۔ خو بيا هم
  $\cfind{e_d} \simeq d$، حال دا چې $d(e_d)$ تعريف شوے دے؛ دغه
  مساوات به بيا $\cfind{e_d}(e_d) \fdefined$ غوښتل۔
\end{enumerate}
د شاخص~$e_d$ په دواړو ممکنو حالتونو کښې تناقض راځي۔ که $h$ جزوي بازګشتي وای، $d$ به هم
جزوي بازګشتي وه او د هغې شاخص~$e_d$ به منل شوی وای؛ خو هېڅ ممکن قيمت
د دغه شاخص له تعريف سره نۀ جوړېږي۔ نو بايد دې پايلې ته ورسېږو چې $h$
جزوي بازګشتي کېدے نۀ شي۔
\end{proof}

\end{document}
